Write \(M_n=2^{dJ_z}\), \(B=B_{J_y,J_z,n}\), \(K=K_{J_y,J_z,n}\), and \(A=A_{J_y,J_z,n}\). All constants below depend only on the density bounds, the compactly supported boundary-wavelet bases, and the fixed support dimension, and hence are uniform over the declared model class.

We first control the raw quadratic kernel without assuming that the two Gram conjugations preserve compact support. Replace the covariate wavelets spanning the levels through \(J_z\) by the orthonormal tensor boundary-scaling basis \(\{\varphi_{J_z,k}:k\in\mathcal K_{J_z}^{\mathrm{sc}}\}\) of the same space. This orthogonal coordinate change leaves the scalar kernel \(\Phi(w)^{\top}B\Phi(w')\), its moments, and the eigenvalues of \(K\) unchanged. Its covariate Gram matrix is
\[
G_z(k,l)=\int_{\mathcal Z_d}\varphi_{J_z,k}(v)\varphi_{J_z,l}(v)r_n(v)\,dv. \tag{1}
\]
By \(\mathrm{ass{:}covariate\text{-}density\text{-}bounds}\),
\[
cI\preceq G_z\preceq CI. \tag{2}
\]
By compact support in \(\mathrm{ass{:}wavelet\text{-}regularity}\), there is a fixed integer \(R\) such that \(G_z(k,l)=0\) whenever the lattice distance between \(k\) and \(l\) exceeds \(R\). Put \(a=(C+c)/2\) and \(E=I-G_z/a\). Equation (2) gives
\[
\|E\|_{\mathrm{op}}\le q:=\frac{C-c}{C+c}<1,
\]
whereas bandedness gives \((E^j)_{kl}=0\) if \(\operatorname{dist}(k,l)>Rj\). The norm-convergent Neumann series therefore yields the entrywise inverse localization
\[
|G_z^{-1}(k,l)|
=\left|a^{-1}\sum_{j\ge0}(E^j)_{kl}\right|
\le \frac{a^{-1}}{1-q}q^{\lceil\operatorname{dist}(k,l)/R\rceil}. \tag{3}
\]
Consequently, for
\[
\kappa_{J_z}(z,v)=\sum_{k,l}\varphi_{J_z,k}(z)G_z^{-1}(k,l)\varphi_{J_z,l}(v),
\]
the wavelet sup-norm bound and bounded overlap imply
\[
|\kappa_{J_z}(z,v)|
\le C_0M_n\exp\{-c_0 2^{J_z}\|z-v\|_2\}. \tag{4}
\]

For the outcome coordinate, let \(R^y_{J_y}\) be the kernel of the orthogonal projector onto the retained zero-mean outcome space. Compact support and bounded overlap give
\[
\sup_{J_y,y}\int_{\mathcal Y}|R^y_{J_y}(y,u)|\,du\le C_1. \tag{5}
\]
The Green kernel of \(\mathcal G_{n,z}\) is uniformly bounded. Indeed, if \(u=\mathcal G_{n,z}h\) and \(\int_0^1h=0\), integration of the differential equation and the Neumann boundary condition gives
\[
u'(y)=-\bar q_n(y,z)^{-1}\int_0^yh(v)\,dv.
\]
The lower density bound in \(\mathrm{ass{:}conditional\text{-}density\text{-}bounds}\), followed by the normalization \(\int_0^1u=0\), represents \(u\) as an integral of \(h\) against a kernel bounded only in terms of \(c\). Combining this representation with (5) on both sides of the inverse shows that the kernel \(g_{J_y}(y,y';z)\) of \(P_{J_y}\mathcal G_{n,z}P_{J_y}\) obeys
\[
\sup_{J_y,y,y',z}|g_{J_y}(y,y';z)|\le C_2. \tag{6}
\]

The full projection Gram matrix is the identity on the orthonormal outcome coordinates tensored with \(G_z\). Expanding both Gram inverses in \(\mathrm{def{:}whitened\text{-}spectral\text{-}array}\) therefore gives the exact identity
\[
b_J\{(y,z),(y',z')\}
:=\Phi(y,z)^{\top}B\Phi(y',z')
=\int_{\mathcal Z_d}\kappa_{J_z}(z,v)\kappa_{J_z}(z',v)
g_{J_y}(y,y';v)\,dR_n(v). \tag{7}
\]
Thus both raw-moment Gram conjugations are included. Equations (4), (6), the upper density bound, and \(\|z-v\|_2+\|z'-v\|_2\ge\|z-z'\|_2\) imply
\[
|b_J\{(y,z),(y',z')\}|
\le C_3M_n\exp\{-c_1 2^{J_z}\|z-z'\|_2\}. \tag{8}
\]
Because each joint arm density is at most \(C^2\), integration of (8) over a covariate tube of volume \(O(M_n^{-1})\) gives, for \(p\in\{2,4\}\),
\[
\max_{a,b}\iint|b_J(w,w')|^p\,dP_{a,n}(w)dP_{b,n}(w')
\le C_4M_n^{p-1},
\]
\[
\max_{a,b,w}\int|b_J(w,w')|^p\,dP_{b,n}(w')
\le C_4M_n^{p-1}. \tag{9}
\]
Let \(\mu_a=E_{P_{a,n}}\Phi(W)\) and
\[
b_{ab}^{\circ}(w,w')=
\{\Phi(w)-\mu_a\}^{\top}B\{\Phi(w')-\mu_b\}.
\]
Expanding the two centerings expresses \(b_{ab}^{\circ}\) as the raw kernel minus its two one-coordinate conditional expectations plus its expectation. Jensen's inequality applied to each conditional expectation preserves the orders in (9). Hence
\[
E|b_{ab}^{\circ}(W_a,W_b)|^2\lesssim M_n,
\qquad
E|b_{ab}^{\circ}(W_a,W_b)|^4\lesssim M_n^3, \tag{10}
\]
and the conditional second moment in either argument has \(L^2\)-norm \(O(M_n)\).

We next record the four-cycle contraction. For
\[
c_{abc}(w,w')=E_{P_{c,n}}
\{b_{ac}^{\circ}(w,W_c)b_{bc}^{\circ}(w',W_c)\},
\]
covariance algebra identifies \(\|c_{abc}\|_2^2\) with a cyclic trace containing two factors \(B\Gamma_cB\) and the two exterior arm covariances. Since \(0\preceq\Gamma_a\preceq\Gamma:=\Gamma_0+\Gamma_1\), each exterior factor can be written as a contraction conjugated by \(\Gamma^{1/2}\). Cyclicity of trace and the Hilbert--Schmidt inequality then give
\[
\|c_{abc}\|_2^2\lesssim\operatorname{tr}(K^4).
\]
By \(\mathrm{prop{:}anisotropic\text{-}spectrum\text{-}bias}\),
\[
\operatorname{tr}(K^4)
\le\|K\|_{\mathrm{op}}^2\operatorname{tr}(K^2)\lesssim M_n,
\qquad
\frac{\operatorname{tr}(K^4)}{\operatorname{tr}(K^2)^2}\lesssim M_n^{-1}. \tag{11}
\]

Index an evaluation observation by \(r=(e,a,i)\), put \(Z_r=\Phi(W_r)-E\Phi(W_r)\), and order observations fold by fold. Direct expansion of the diagonal-free statistic in \(\mathrm{def{:}anisotropic\text{-}estimator}\), using \(n=2m\), gives
\[
h_{n,rs}(x,y)=
\begin{cases}
\dfrac{2}{m-1}x^{\top}By,&e_r=e_s,\ a_r=a_s,\\[4pt]
-\dfrac{2}{m}x^{\top}By,&e_r=e_s,\ a_r\ne a_s,\\[4pt]
0,&e_r\ne e_s,
\end{cases}
\qquad
n\mathbb U_{J_y,J_z,n}=\sum_{r<s}h_{n,rs}(Z_r,Z_s). \tag{12}
\]
Independence in \(\mathrm{ass{:}two\text{-}arm\text{-}sampling}\) and \(EZ_r=0\) show that each kernel is canonical in either argument. There are \(O(m)\) nonzero edges incident to a vertex and \(O(m^2)\) nonzero edges in total. Equations (10) and (12) therefore give
\[
\max_r\sum_{s\ne r}Eh_{n,rs}^2\lesssim\frac{M_n}{m},
\qquad
\sum_{r<s}Eh_{n,rs}^4\lesssim\frac{M_n^3}{m^2}. \tag{13}
\]
Canonicality makes distinct edges orthogonal, including two edges sharing exactly one vertex. Summing the edge variances and separating within-arm from cross-arm edges yields
\[
\begin{split}
\operatorname{Var}(n\mathbb U_{J_y,J_z,n})
&=\frac{4m}{m-1}\sum_{a=0}^1\operatorname{tr}(B\Gamma_aB\Gamma_a)
+8\operatorname{tr}(B\Gamma_0B\Gamma_1)\\
&=4\operatorname{tr}\{B(\Gamma_0+\Gamma_1)B(\Gamma_0+\Gamma_1)\}
+O(m^{-1})\operatorname{tr}(K^2)\\
&=A^2\{1+O(m^{-1})\}.
\end{split} \tag{14}
\]
Here \(A^2=4\operatorname{tr}(K^2)\asymp M_n\) by \(\mathrm{prop{:}anisotropic\text{-}spectrum\text{-}bias}\), proving the variance assertion.

It remains to verify Lindeberg and predictable-variance concentration. Let
\[
q_r=\sum_{s<r}h_{n,sr}(Z_s,Z_r),
\qquad
\mathcal D_{n,r}=\ell_{n,r}(W_r)+q_r.
\]
Canonicality gives \(E(\mathcal D_{n,r}\mid\mathcal F_{r-1})=0\) and \(\sum_r\mathcal D_{n,r}=G_n+n\mathbb U_{J_y,J_z,n}\). Conditional on \(W_r\), the summands in \(q_r\) are independent and centered. Thus
\[
Eq_r^4
=\sum_{s<r}Eh_{n,sr}^4
+6\sum_{s<t<r}E\!\left[E(h_{n,sr}^2\mid W_r)E(h_{n,tr}^2\mid W_r)\right]. \tag{15}
\]
The sectional bound after (10), including the \(O(m^{-1})\) edge coefficients, bounds the sum of conditional second moments in (15) by \(C_5M_n/m\). Summing over \(r\) and applying (13) gives
\[
\sum_rEq_r^4\lesssim\frac{M_n^2}{m}+\frac{M_n^3}{m^2}. \tag{16}
\]
If \(|\ell_{n,r}|\le L_0\), independence and centering imply
\[
\sum_rE\ell_{n,r}^4\le L_0^2\sum_rE\ell_{n,r}^2
=L_0^2EG_n^2=O(S_n^2).
\]
Since \(S_n\ge A\asymp M_n^{1/2}\), (16) and \((x+y)^4\le8(x^4+y^4)\) yield
\[
S_n^{-4}\sum_rE\mathcal D_{n,r}^4
\lesssim M_n^{-1}+m^{-1}+\frac{M_n}{m^2}\longrightarrow0. \tag{17}
\]
Indeed, \(J_z\to\infty\) gives \(M_n\to\infty\), while \(M_n\le D_{J_y,J_z}=o(n^2)=o(m^2)\). For every \(\epsilon>0\), the bound
\[
x^2\mathbf 1\{|x|>\epsilon S_n\}\le\frac{x^4}{\epsilon^2S_n^2}
\]
and (17) prove the stated Lindeberg conclusion after summation and expectation.

Finally, direct expansion gives
\[
\sum_rE(\mathcal D_{n,r}^2\mid\mathcal F_{r-1})
=E(G_n+n\mathbb U_{J_y,J_z,n})^2+C_n+R_n+T_n, \tag{18}
\]
where
\[
\begin{split}
C_n&=2\sum_r\sum_{s<r}E\{\ell_{n,r}(W_r)h_{n,sr}(Z_s,Z_r)\mid Z_s\},\\
R_n&=\sum_r\sum_{s<r}\{E(h_{n,sr}^2\mid Z_s)-Eh_{n,sr}^2\},\\
T_n&=2\sum_r\sum_{s<t<r}E\{h_{n,sr}h_{n,tr}\mid Z_s,Z_t\}.
\end{split} \tag{19}
\]
To verify the constant in (18), canonicality gives \(E(G_n\,n\mathbb U_{J_y,J_z,n})=0\), and the orthogonality of distinct canonical edges gives \(E(n\mathbb U_{J_y,J_z,n})^2=\sum_{r<s}Eh_{n,rs}^2\). Thus the constant terms from the conditional squares are exactly \(E(G_n+n\mathbb U_{J_y,J_z,n})^2\).

The repeated-edge term \(R_n\) groups into independent centered functions of \(Z_s\). The conditional-second-moment \(L^2\) bound after (10), the \(O(m^{-1})\) edge coefficients, and the \(O(m)\) future edges give
\[
ER_n^2\lesssim\frac{M_n^2}{m}. \tag{20}
\]
For \(T_n\), regroup first over the unordered pair \((s,t)\). Pair kernels with disjoint edge sets are independent, and those sharing only one vertex are orthogonal by canonicality. A surviving repeated pair is the four-cycle contraction in (11); summation over possible common terminal vertices has total coefficient \(O(m^{-1})\). Since there are \(O(m^2)\) pairs,
\[
ET_n^2\lesssim m^2m^{-2}\operatorname{tr}(K^4)\lesssim M_n. \tag{21}
\]
For the mixed contraction, put \(u_r=E\{Z_r\ell_{n,r}(W_r)\}\). Covariance Cauchy--Schwarz gives
\[
\|\Gamma_{a_r}^{-1/2}u_r\|_2^2\le E\ell_{n,r}^2,
\]
with the inverse understood on the covariance range. Moreover,
\[
\|\Gamma_a^{1/2}B\Gamma_b^{1/2}\|_{\mathrm{op}}
\le\|\Gamma^{1/2}B\Gamma^{1/2}\|_{\mathrm{op}}
=\|K\|_{\mathrm{op}}\lesssim1,
\]
because \(0\preceq\Gamma_a,\Gamma_b\preceq\Gamma\). Regrouping \(C_n\) over \(s\), applying Cauchy--Schwarz to its \(O(m)\) future indices, and then using independence of the resulting centered functions of \(Z_s\) proves
\[
EC_n^2\lesssim EG_n^2=O(S_n^2). \tag{22}
\]
Equations (20)--(22), \((x+y+z)^2\le3(x^2+y^2+z^2)\), \(S_n^4\ge A^4\asymp M_n^2\), and \(S_n^2\gtrsim M_n\) now give
\[
\frac{E\left[\left\{\sum_rE(\mathcal D_{n,r}^2\mid\mathcal F_{r-1})
-E(G_n+n\mathbb U_{J_y,J_z,n})^2\right\}^2\right]}{S_n^4}
\lesssim m^{-1}+\frac{M_n}{m^2}+M_n^{-1}+o(1)\longrightarrow0. \tag{23}
\]
Equations (1)--(23) establish the claimed noncompact Gram-inverse localization, both conjugations, fourth moments, maximal influence, repeated-edge, shared-vertex, four-cycle and mixed contractions, Lindeberg condition, and predictable-variance concentration.